% Generated draft: requires mathematical and source review.
\documentclass[11pt]{article}
\usepackage[margin=25mm]{geometry}
\usepackage{fontspec}
\setmainfont{Latin Modern Roman}
\usepackage{amsmath,amssymb,mathtools,unicode-math}
\setmathfont{Latin Modern Math}
\usepackage{enumitem,needspace,xurl,hyperref}
\setlength{\parindent}{0pt}
\setlength{\parskip}{5pt}
\setlength{\emergencystretch}{3em}
\newcommand{\N}{\mathbb{N}}
\newcommand{\Z}{\mathbb{Z}}
\newcommand{\Q}{\mathbb{Q}}
\newcommand{\R}{\mathbb{R}}
\newcommand{\C}{\mathbb{C}}
\newcommand{\F}{\mathbb{F}}
\newcommand{\A}{\mathbb{A}}
\newcommand{\PP}{\mathbb{P}}
\newcommand{\E}{\mathbb{E}}
\newcommand{\eps}{\varepsilon}
\newcommand{\dd}{\,\mathrm{d}}
\begin{document}
\section*{UnsolvedMath ID: 30001794}
\textbf{Source problem code:} catalogue entry
\subsection{Definitions and mathematical statement}
The supplied question asks whether all universal sofic groups \(S_U\) are elementarily equivalent. The record does not define \(S_U\) or specify the logical language. The following gives a conventional construction and a precise conditional formulation; its identification with the intended source formulation requires review.

For each integer \(n\geq 1\), let \(\operatorname{Sym}(n)\) be the group of permutations of \(\{1,\ldots,n\}\), with normalized Hamming distance
\[
 d_n(\sigma,\tau)=\frac{1}{n}\bigl|\{i\in\{1,\ldots,n\}:\sigma(i)\ne\tau(i)\}\bigr|.
\]
Let \(U\) be a nonprincipal ultrafilter on the positive integers, meaning an ultrafilter containing no finite set. Define the normal subgroup
\[
 N_U=\left\{(\sigma_n)_{n\geq1}\in\prod_{n\geq1}\operatorname{Sym}(n):\lim_{n\to U}d_n(\sigma_n,1)=0\right\}
\]
and the metric ultraproduct
\[
 S_U=\left(\prod_{n\geq1}\operatorname{Sym}(n)\right)/N_U,
 \qquad
 d_U([\sigma_n],[\tau_n])=\lim_{n\to U}d_n(\sigma_n,\tau_n).
\]
Here \(\lim_{n\to U}\) denotes the ultrafilter limit of a bounded real sequence.

In the ordinary first-order language of groups \(L_{\mathrm{grp}}=\{\cdot,{}^{-1},1\}\), two groups are elementarily equivalent if they satisfy exactly the same first-order sentences without parameters. Under the construction above and this language convention, the question is whether
\[
 \text{for every pair of nonprincipal ultrafilters }U,V\text{ on the positive integers,}
 \quad S_U\equiv_{L_{\mathrm{grp}}}S_V.
\]
Equivalently, for every such \(U,V\) and every \(L_{\mathrm{grp}}\)-sentence \(\varphi\), does
\[
 S_U\models\varphi\quad\Longleftrightarrow\quad S_V\models\varphi
\]
hold?

If elementary equivalence is instead intended in continuous logic for groups equipped with their normalized Hamming metric, the corresponding question is whether every continuous sentence in that metric group language has the same value in \((S_U,d_U)\) and \((S_V,d_V)\), for every such \(U,V\). These are distinct interpretations of the question; the supplied record does not resolve the choice. It also does not specify whether other sequences of permutation degrees are included in the quantifier over universal sofic groups.
\subsection{Short English statement}
Do all universal sofic groups satisfy the same logical sentences? Using the conventional ultraproduct of finite symmetric groups, this asks whether changing the nonprincipal ultrafilter changes the elementary theory. The source record leaves the construction conventions and the choice between ordinary group logic and continuous metric logic unspecified.
\subsection{Sources}
\begin{itemize}
\item Local UnsolvedMath registry, permanent ID \(30001794\), title \emph{Elementary Equivalence of Universal Sofic Groups}, source problem code \texttt{catalogue entry}. The supplied record gives proposed year \(2011\) and names no proposer.
\item \emph{Finite-dimensional Approximations of Discrete Groups} (2011), Oberwolfach Reports, pp.\ 1429--1467, DOI \texttt{10.4171/owr/2011/26}: \url{https://doi.org/10.4171/owr/2011/26}. The supplied background describes Oberwolfach Reports as documenting mathematical workshops and their open-problem discussions; no problem-specific page is supplied.
\item Reference labeled ``Thomas'' in the supplied literature sources: \url{https://sites.math.rutgers.edu/~sthomas/Universal-Sofic.pdf}. The record supplies no full bibliographic citation for this reference.
\item Supplied literature assessment, checked on \texttt{2026-08-21}: status ``open''; the question is described as unresolved, with existing work giving only partial theory comparisons. This is supplied dated curation derived from the OWR record, not an independently verified assessment. The record expressly requests independent verification of time-sensitive status and cited claims.
\item The registry labels the statement \texttt{corrected\_verified} and reports manual semantic, typographic, and LaTeX review of its cleanup. This is supplied editorial metadata, not independent verification of the construction, logical interpretation, or mathematical status.
\end{itemize}
\end{document}
